\documentclass[11pt]{article}
\usepackage{amsmath,amssymb,amsthm,hyperref}
\newtheorem{theorem}{Theorem}\newtheorem{proposition}[theorem]{Proposition}
\newcommand{\B}{\mathcal B}\newcommand{\E}{\mathcal E}\newcommand{\D}{\mathcal D}
\title{Native homogeneous polynomial substitution with shared coupling and exact state closure}
\date{16 November 2026}
\begin{document}\maketitle
\paragraph{Status.} The full written argument has passed independent review, including the native coefficient and content maps, substitution coherence, and positive-degree coreflection. No proof-assistant certification or TeX compilation is claimed. This is a genuine multi-occurrence extension of the native coupling-span bridge, with a declared object-level boundary: every global content grade has one fixed finite arity, though its shape and all output values may vary with the actual preceding witness. Grades containing shapes of different arities are excluded. A variable-arity compiler is not established here; no unrestricted O1 or v0.2 completion is claimed.

\section{Independent original inputs}
Let $\mathcal A$ be a locally small symmetric rig whose additive tensor is categorical finite coproduct, with canonical distributors, and assume that $\E=\mathcal A/I$ has finite limits. The original O1 rig
\[
 \mathcal V=\mathbb R_{\ge0}^{\mathrm{disc}}\times\mathcal A
\]
has finite augmented families exactly equivalent to the native category $\mathcal C$ of finite raw probability spaces $(W,p)$ with one coefficient $C_w\in\E$ at each raw point. Native arrows are actual pushforward-preserving functions and actual coefficient arrows over the augmentation. Indeed the scalar augmentation says $\sum_wp_w=1$, and the coproduct universal property decomposes the coefficient augmentation and every fibre aggregation into their individual arrows. Canonical distributivity identifies the old dependent and independent coefficient operations with the indicated tensors. This equivalence retains all native maps, not just a law shadow. Finite-dimensional vector spaces with direct sum and tensor give a nonfaithful coefficient example.

Independently choose a locally small symmetric monoidal category $\D$ for content coefficients. A homogeneous global content grade is
\[
 C=(n,O,(D_o)_{o\in O}),\quad n\in\mathbb N,
\]
where $O$ is a nonempty small set of shapes, each with the ordered position set $[n]$. A grade map has the same arity, a function $\alpha:O\to O'$, and actual arrows $D_o\to D'_{\alpha(o)}$. Its extension at a set $B$ is $O\times B^n$; it forgets the coefficient objects but not the shape. This is a constant O1 content category, with all same-arity shape merges and all coefficient arrows. The nonempty-shape convention avoids assigning several different degrees to the same empty polynomial collection; it is an explicit further object restriction.

Executing a grade $C=(n,O,D)$ and then a grade $E=(m,P,F)$ at every occurrence has global grade
\[
 C\circ E=(nm,O\times P^n,\ D_o\otimes\bigotimes_{i=1}^n F_{p_i}). \tag{G}
\]
In the original contentad's execution convention this is $E\diamond C$, since its extension is $F_C(F_E B)$. The unit is the unary single-shape unit coefficient. Associativity is the full shape-tuple bijection and the canonical permutation of tagged coefficient factors; it does not delete unused shapes or identify different coefficient arrows. Empty position products when $n=0$ have their original nullary meaning. Thus an entirely unexecuted continuation grade contributes zero factors, just as in the original polynomial substitution.

This content input satisfies the original constant-fibration display and Beck--Chevalley conditions. Its shape extension is cartesian under every pure map, since each shape and ordered position set are fixed; arity is the ordered-list map. Native grade maps preserve that arity. Unit and substitution act on all shape/coefficient components, with their coherence given by the same addressed tree. No mixed interaction is postulated as an input.

\section{Target joint-substitution bicategory}
Let $\B=\operatorname{Fam}(\mathcal C)$, with all small index families. The native multilimit theorem gives finite limits in $\B$. Let $U:\B\to\mathbf{Set}$ be total raw carrier; no finite-limit preservation by $U$ is assumed.

Define a horizontal process $P:X\to Y$ to consist of its entire global grade $C=(n,O,D)$, an apex $Z\in\B$, a left leg $s:Z\to X$, $n$ right legs $t_i:Z\to Y$, and a label function $\ell:UZ\to O$. A 2-cell has the same degree $n$, an actual apex map commuting with the legs, and one full global grade map satisfying the label equation. All unused global shapes and coefficients remain part of this 2-cell, even when the apex is empty.

For $P:X\to Y$ of degree $n$ with apex $Z$, and $Q:Y\to V$ of degree $m$ with apex $W$, set
\[
 Z_{QP}=Z\times_{Y^n}W^n. \tag{C}
\]
Its left leg is the old left leg through $Z$. Its $nm$ right legs are indexed by $(i,j)\in[n]\times[m]$ and use the $j$th output of the $i$th $W$ projection. Its global grade is (G), and its label pairs the root label with the $n$ continuation labels through those projections. The identity has degree one, apex $X$, both legs identity, and unit global grade.

\begin{theorem}[Actual polynomial substitution]
These data form a bicategory with native maps as 2-cells. Composition is genuinely dependent on the actual output values through (C), permits arbitrary joint couplings across its occurrence apices, and retains actual coefficient and global-grade arrows. The construction also has changing-pure-endpoint squares, obtained from arbitrary maps of objects of $\B$ and the corresponding leg equations.
\end{theorem}
\begin{proof}
For three degrees $n,m,k$, the two apex bracketings are
\[
 (Z\times_{Y^n}W^n)\times_{V^{nm}}T^{nm},\qquad
 Z\times_{Y^n}(W\times_{V^m}T^m)^n.
\]
Both are limits of the same diagram with one $Z$, $n$ copies of $W$, and $nm$ copies of $T$, and precisely the same input/output equations. At family-index level, flattening an iterated coupling pushes its law to the complete addressed tuple of all these primitive witnesses. Both parenthesizations retain that same complete joint law, not merely its primitive marginals; conversely its projections determine the intermediate coupling indices used in either factorization. Finite powers preserve pullbacks, so their canonical comparison is invertible. The global shape comparison is
\[
 (O\times P^n)\times R^{nm}\cong O\times(P\times R^m)^n.
\]
On coefficients it is the symmetry/associativity map with the same tagged root, child and grandchild factors. On labels it is exactly that tuple comparison. Thus it is a full native process isomorphism. Both routes around the pentagon project to the same four-level primitive diagram; limit uniqueness proves the apex equation, and symmetric monoidal coherence on distinctly tagged factors proves the coefficient equation. Unit equations use singleton powers, pullback identities, and the actual grade unitors. Native 2-cell composition uses induced pullback maps and actual global-grade substitution; their naturality follows from their projection equations. This argument neither replaces native arrows by isomorphism classes nor assumes that $U$ preserves the displayed limits. For $n=0$, the empty powers are terminal and both formulas reduce to the untouched root apex, agreeing with nullary content substitution.
\end{proof}

\begin{proposition}[Native universal characterization of one substitution]
For every $T\in\mathcal B$, maps $T\to Z_{QP}$ are naturally in bijection with a map $T\to Z$ and $n$ maps $T\to W$ whose intermediate endpoint maps agree at each ordered occurrence. Thus the apex retains every admitted joint law and every actual compatible coefficient cone; it selects no universal probability coupling. On each component, its raw witnesses are the full matching tuples, its family indices are exactly their joint laws with the stated marginals, and its coefficients are the corresponding native coefficient limits. The composite content grade and body are determined by their displayed full tuple formulas.
\end{proposition}
\begin{proof}
This is the pullback/product universal property in $\mathcal B$, evaluated using the native multilimit formula. Each native cone forces its own joint-law index by pushforward; coefficient limit uniqueness supplies its actual factorization. Naturality follows by precomposition. It does not assert that one exact-law object represents all such cones in the uncompleted native category.
\end{proof}

No symmetric monoidal bicategory structure on these varying-degree homs is asserted merely from products of objects. Two degrees need not combine into the same output type under an external tensor. The theorem concerns the displayed polynomial substitution operation, not an unproved O2 tensor.

\section{The full original-source bridge}
Restrict original interfaces to finite sets. A source process of global degree $n$ has the original native row data $(W_a,\mu_a,C_{a,w})$, its entire global grade $C$, and body
\[
 w\longmapsto\big(o(a,w),b_1(a,w),\ldots,b_n(a,w)\big).
\]
Its composition with degree-$m$ data has dependent witnesses $(w,z_1,\ldots,z_n)$, law
\[
 \mu_a(w)\prod_i\nu_{b_i(a,w)}(z_i),
\]
context coefficient $C_{a,w}\otimes\bigotimes_i E_{b_i,z_i}$, and full substituted global grade (G). These are precisely the original independent O1 crossing and mixed composition for the independently specified inputs above, not a unary encoding.

For a finite set $A$, let $LA$ be the all-prior family $(A,p,I)_{p\in\mathsf D A}$. Map a source process to the target apex $\Gamma P$ whose component at $p$ has raw set $\coprod_a W_a$, law $p(a)\mu_a(w)$, coefficient $C_{a,w}$, source map $a$, output maps $b_i$, and global shape label $o$. The original entire global grade and every native grade arrow are retained unchanged.

\begin{theorem}[Locally faithful oplax polynomial bridge]
This gives a locally faithful oplax functor with invertible unit constraints from the original homogeneous O1 mixed process bicategory to the target, with directed compositor
\[
 L(QP)\longrightarrow LQ\circ LP.
\]
It preserves changing-pure-endpoint squares. Its compositor is not generally invertible, and the target has additional prior-indexed native cells.
\end{theorem}
\begin{proof}
At prior $p$, choose the complete law
\[
 \lambda_p(a,w,\mathbf z)=p(a)\mu_a(w)\prod_i\nu_{b_i(a,w)}(z_i)
\]
on the pullback raw witnesses of (C). Its root marginal is the $\Gamma P$ law and its $i$th continuation marginal is the $\Gamma Q$ law at the pushed-forward prior of $b_i$. Hence it selects a component of the native pullback family. Map the original dependent tuple to that component by its canonical coordinate bijection and use the augmentation-induced coefficient arrow
\[
 C_{a,w}\otimes\bigotimes_i E_{b_i,z_i}
 \longrightarrow C_{a,w}\times\prod_i E_{b_i,z_i}.
\]
The global grade map is identity on the full (G), including unused shape combinations and all their coefficients. The label and output equations are the actual original body equations.

For triple composition, both comparison routes select the same complete law, the product of conditional probabilities over the full addressed tree. Both coefficient arrows have exactly the same projection to each primitive factor; augmentation and symmetric monoidal coherence, followed by product uniqueness, prove equality of actual arrows. The grade maps are the original addressed substitution associators. This proves the oplax coherence, not just equality of observable laws.

A source native cell induces the same witness and coefficient maps at every prior and the same single global grade map. Equality of images implies equality of each original map, since every raw point, including zero-weight points, remains in every prior component. If the input is empty, the global grade and its arrows are still retained explicitly, so faithfulness remains true. Pure input maps push priors forward and send $(p,a,w)$ to $(f_*p,f(a),r_a(w))$; the native law equation follows by summing $p(a)\mu_a(w)$ over each fibre. Output-coordinate and label squares are those supplied by the original native square. This also proves the endpoint naturality of the compositor. No inverse native comparison has been introduced.
\end{proof}

A binary constant-output shape followed by two fair-bit continuations separates the compositor: the original law makes the two continuation witnesses independent, whereas (C) also has diagonal and antidiagonal couplings. With augmented vector-space coefficients this remains a nonthin native example; the coefficient maps are not inferred from the raw separation.

The actual coefficient comparison can also be strictly noninvertible independently of the effect choice. In $\E=\operatorname{FinVect}_k/k$, take two continuation coefficients $k^2$, each augmented by its first coordinate. Their independent tensor has dimension four; their categorical product in $\E$ is the pullback $k^2\times_k k^2$, of dimension three. With basis $e_0,e_1$ and augmentation $\epsilon(e_0)=1,\epsilon(e_1)=0$, the canonical comparison sends
\[
 e_0\otimes e_0\mapsto(e_0,e_0),\quad
 e_1\otimes e_0\mapsto(e_1,0),\quad
 e_0\otimes e_1\mapsto(0,e_1),\quad
 e_1\otimes e_1\mapsto0.
\]
Thus its nontrivial kernel is a native grade distinction, not an artifact of forgetting witnesses. Local faithfulness of the bridge retains the original objects and arrows before this comparison; it does not turn the oplax compositor into a faithful or invertible operation-preservation map.

\section{Shared-state closure is not an imported unary trace}
Take a target process $P:X\times LS\to Y\times LS$ of degree $n\ge1$, where $S$ is a finite pure state set. Write $s_S$ for its input-state leg and $t_{i,S}$ for its $i$th output-state leg. Define shared-state closure by the simultaneous equalizer
\[
 Z^{\mathrm{cl}}\longrightarrow Z,
 \qquad s_S=t_{i,S}\quad\text{for every }i\in[n]. \tag{F}
\]
The outside input and $n$ outside outputs restrict through this equalizer. The full global content grade and label are retained, the latter by restriction only. This produces a degree-$n$ process $X\to Y$. For degree zero, deleting an unused state input is a different hiding operation; it is not called feedback here.

\begin{theorem}[Derived feedback laws and sequential comparison]
Shared-state closure is functorial on native cells and pure external endpoint squares. A compatible pure state map, using the same map on the input and on every output state leg, induces a directed native closure comparison. Closing two finite shared-state coordinates in either order is canonically isomorphic to closing their product, with the corresponding nesting coherence. Write $\operatorname{rcl}(P)$ for the same equality restriction while retaining the original state-bearing endpoints $X\times LS$ and $Y\times LS$. For positive degrees $n,m$, there is a canonical directed comparison
\[
 \operatorname{rcl}(Q)\operatorname{rcl}(P)
 \longrightarrow \operatorname{rcl}(QP). \tag{S}
\]
Hiding the common state only after this assembly gives a comparison to $\operatorname{cl}(QP)$. Composing separately hidden closures is a different operation and has no such general projection-preserving comparison. These comparisons are coherent under further substitution but are not generally invertible.
\end{theorem}
\begin{proof}
A native cell commuting with all input/output state legs carries every equality in (F) to the target equality. Equalizer universality constructs the actual map. For a pure state map $h$, source equality implies equality after applying $Lh$; reflection is not needed for this directed map. Target closure may admit extra components, so invertibility is not asserted. External endpoint squares restrict unchanged.

For two state coordinates, either order imposes precisely all $S$ and all $T$ equalities on the same apex. Their limit comparisons are canonical isomorphisms, and all nesting diagrams commute by uniqueness on that full finite diagram. Global grade maps remain identity up to the declared substitution/unit constraints.

For (S), the retained endpoints glue the state as well as the external intermediate value. Its left apex therefore retains a root state $s$, each intermediate state $s_i$, and final states $s_{ij}$. Its equations include $s_i=s$ and $s_{ij}=s_i$. Therefore it satisfies every equation $s_{ij}=s$ required on the right. Projection into the unclosed joint apex followed by equalizer universality gives (S). On actual context coefficients the map is the canonical limit map, and on global content grades it is the same full substituted grade. For three stages both routes impose the same stronger intermediate equations and map to the same final equations; uniqueness proves coherence. No serializability or independence of shared couplings is assumed.
\end{proof}

\begin{theorem}[Homwise coreflection and lax assembly]
Fix a finite pure state $S$ and restrict the ambient state-bearing hom categories to positive-degree processes. In each of these restricted hom categories, let the closed subcategory be the full subcategory whose input-state leg and every output-state leg already agree. The positive-degree ambient restriction is closed under composition and contains identities; degree-zero processes remain in the unrestricted substitution bicategory but are outside this coreflection theorem. Its inclusion $J$ has $\operatorname{rcl}$ as a right adjoint. The counit $J\operatorname{rcl}(P)\to P$ is the native equalizer inclusion with identity global grade map. These closed subcategories are closed under horizontal composition and contain identities. Their coreflections therefore assemble with the canonical lax compositor (S) and invertible unit comparison. This is a homwise universal characterization of this shared-state closure, not a free universal characterization of the whole native extension.
\end{theorem}
\begin{proof}
Let $K$ already have equal state legs and let $(h,\beta):K\to P$ be a native process cell. Its apex map $h$ equalizes every pair defining (F), because the corresponding source legs agree and the cell respects the state-bearing endpoints. Therefore $h$ factors uniquely through the equalizer apex. Retain exactly the same full global grade map $\beta$; its label equation restricts along the factorization. This gives a natural bijection between cells $K\to J\operatorname{rcl}(P)$ and cells $JK\to P$. No body labels or unused grade coefficients are discarded by the factorization. This proves the adjunction.

If $K$ and $L$ already have equal state legs, their compositional pullback matches the output state of each $K$ occurrence to the input state of its $L$ continuation. Consequently all final states equal the common initial state. Their composite is closed; identities are closed as well. Compose the two coreflection counits and apply the proved universal property to factor through $\operatorname{rcl}(QP)$. This is exactly (S). The unit is the canonical isomorphism for an already closed identity. Every associativity/unit coherence diagram has the same composite to the unclosed process; uniqueness of the coreflection factorization proves it. The right adjoints are not claimed to form a pseudofunctor: their compositor need not be invertible, as the double-toggle example demonstrates.
\end{proof}

For pure state objects $LS$, the equalizer description is explicit. Each surviving joint-law component must already be supported on $s=t_{i,S}$ for every output port; then restriction has total mass one. Its raw equality subset retains every closed zero-weight witness and may remove nonclosed zero-weight witnesses. Context coefficients are unchanged because all coefficient legs to the pure state have terminal targets. This is not conditioning and not merely equality of state marginals. A toggling state kernel has a stationary fair law but no fixed-loop component.

Separately hiding state before composition loses the required matching equation: even for state identities, the two hidden factors can retain different state values, and no map to a common-state composite can preserve both state projections. The retained-endpoint formulation of (S) is therefore essential.

The comparison (S) is strictly directed even for deterministic branching: let $S$ be a bit, let $P$ emit $n$ copies of its flipped state, and let every occurrence of $Q$ emit $m$ copies of its flipped input state. Each separate closure is empty. The composite returns the original state at all $nm$ outputs, so its joint closure is nonempty. Thus independently closing branches is not identified with closing their assembled shared interaction. The degree-one case recovers the earlier span fixed-loop closure; the higher-degree law was derived from its multiple equality constraints rather than assumed to be an ordinary unary trace.

\section{Scope and unresolved variable arity}
The result supplies an actual original-O1 source class, a coherent native joint multi-occurrence composition, its locally faithful directed source bridge, and explicit shared-state closure with nesting and substitution comparison. Shape merges and coefficient arrows within a fixed arity remain unrestricted; output values and continuation rows depend on actual preceding witnesses; global unused shapes and coefficients are not dropped.

The global-object restriction is material. A grade containing both unary and binary shapes is not homogeneous, even if every native grade map separately preserves each shape's arity. Variable arity can produce outcome-dependent position reassignment under flattening. The present proof uses static addressed index arrays and does not establish stability of arbitrary shared-law admission under such reassignment or under native maps forgetting ancestry. Merely retaining a history is not a proof of that missing stability. A wider completion needs a proved routing-stable admission/comprehension law or a different declared comparison direction. Infinite recurrence, infinite-effect completion, and a universal characterization of the whole original extension remain open. This theorem is a substantive compatible relative construction, not a relabelling of the general v0.2 target as solved.
\begin{thebibliography}{9}
\bibitem{O1} \emph{General Independent Interaction Criterion}, project source, 1 November 2026. The original augmented-rig native arrows, constant polynomial content and mixed composition.
\bibitem{Multi} \emph{Native multilimits and the dependent-selector obstruction}, project source, 16 November 2026. The finite-limit construction in all-small-family native effects.
\bibitem{Bridge} \emph{A directed bridge from conditional native bind to all-coupling span feedback}, including its global-writer lift, project source. The degree-one predecessor and its declared preservation boundaries.
\end{thebibliography}
\end{document}
